\documentclass[11pt]{article}
\usepackage[margin=1in,top=0.9in,bottom=0.9in]{geometry}
\usepackage{amsmath,amssymb,amsthm,hyperref}
\hypersetup{colorlinks=true,linkcolor=blue,urlcolor=blue,citecolor=blue}
\setlength{\emergencystretch}{2em}
\newtheorem{theorem}{Theorem}
\newtheorem{lemma}[theorem]{Lemma}
\newcommand{\VS}{\operatorname{VS}}
\newcommand{\LCR}{\operatorname{D}}
\title{Negative-order Vendi--LCR comparison\\without a full-rank assumption}
\author{Research proof and formalization packet}
\date{October 26, 2026}
\begin{document}
\maketitle
\begin{abstract}
For every finite positive semidefinite similarity matrix, we prove that
the Vendi score bounds the uniform Leinster--Cobbold--Reeve diversity at
all negative orders, including the directly defined minus-infinity
endpoint. Zero eigenvalues are omitted, following the source's archived
implementation. A weighted spectral Jensen inequality gives a stronger
bound through the diagonal of the matrix square. The proof requires no
full-rank assumption, and the complete matrix-to-score argument is
verified in Lean. Known nonnegative-order and full-rank results are
credited; classical Jensen and H\"older--McCarthy inequalities are
background ingredients. Source-access and prior-art limits are stated
explicitly.
\end{abstract}

\section{Definitions and the remaining source case}
Let $Z$ be a real symmetric positive semidefinite $N\times N$ matrix,
where $N\geq1$, with
\[
 Z_{ii}=1,\qquad 0\leq Z_{ij}\leq1.
\]
Write $\lambda_1,\ldots,\lambda_N\geq0$ for its eigenvalues, counted
with multiplicity, and put
\[
 J=\{j:\lambda_j>0\},\qquad R_i=\sum_j Z_{ij},\qquad
 E_i=(Z^2)_{ii}=\sum_j Z_{ij}^2.
\]
Thus $1\leq E_i\leq R_i$, and $J$ is nonempty. For finite real $q<1$
define the uniform scores by
\begin{equation}\label{eq:scores}
 \VS_q(Z)=\left(\sum_{j\in J}(\lambda_j/N)^q\right)^{1/(1-q)},
 \qquad
 \LCR_q(Z)=\left(\frac1N\sum_i(R_i/N)^{q-1}\right)^{1/(1-q)}.
\end{equation}
The positive-spectrum index set retains eigenvalue multiplicities.
For $q<0$, no negative power is assigned to a zero eigenvalue.

This convention is material. The general displayed spectral formula in
Nguyen et al.\ \cite{NguyenPreprint} does not separately explain singular
negative powers, but its Appendix~C.3 explicitly filters out zero
eigenvalues before taking powers, including in searches with negative
orders and low-rank Gram matrices. We use that finite support formula.
The alternative extended-real convention $0^q=+\infty$ would give a
trivial infinite singular comparison; full-rank regularization toward
a singular matrix would not prove \eqref{eq:scores}.

Theorem~1 of the latest listed arXiv revision, v4 of March~30, 2026,
proves $\VS_q\geq\LCR_q$ for all $q\geq0$, and for $q<0$ when $Z$ is
full rank \cite{NguyenPreprint}. The public abstract of the May~26, 2026
journal version retains that restriction \cite{NguyenJournal}.
The result below supplies the arbitrary-rank negative-order case of
the comparison represented by MathDB~371345 and its equivalent
power-sum formulation, MathDB~371346. These are formulations of one
comparison. No novelty is claimed for the already proved cases.

\section{A spectral energy bound}
\begin{theorem}\label{thm:energy}
For every matrix $Z$ as above and every real $q<1$,
\begin{equation}\label{eq:energy}
 \sum_{j\in J}\lambda_j^q
 \ \geq\ \sum_i E_i^{q-1}
 \ \geq\ \sum_i R_i^{q-1}.
\end{equation}
Consequently $\VS_q(Z)\geq\LCR_q(Z)$ for every finite real $q<1$;
in particular this holds for every $q<0$ at arbitrary rank.
\end{theorem}
\begin{proof}
Choose a real orthonormal eigenbasis $v_1,\ldots,v_N$, and set
$u_{ij}=v_j(i)^2$. Spectral decomposition and column orthonormality give
\begin{equation}\label{eq:frames}
 \sum_{j\in J}\lambda_j u_{ij}=1,\qquad
 \sum_{j\in J}\lambda_j^2u_{ij}=E_i,\qquad
 \sum_i u_{ij}=1.
\end{equation}
The omitted zero-eigenvalue terms vanish in the first two identities.
For each $i$, the nonnegative numbers $w_{ij}=\lambda_j u_{ij}$,
$j\in J$, therefore sum to one, with weighted mean of the positive
eigenvalues equal to $E_i$.

The function $t\mapsto t^{q-1}$ is convex on $(0,\infty)$ when $q<1$.
Weighted scalar Jensen, followed by summation over $i$, gives
\[
 E_i^{q-1}\leq\sum_{j\in J}w_{ij}\lambda_j^{q-1}
       =\sum_{j\in J}\lambda_j^q u_{ij},
 \qquad
 \sum_i E_i^{q-1}\leq\sum_{j\in J}\lambda_j^q.
\]
Since $E_i\leq R_i$ and $q-1<0$, we also have
$E_i^{q-1}\geq R_i^{q-1}$, proving \eqref{eq:energy}.

Both normalized interiors in \eqref{eq:scores} have the same factor:
\begin{align*}
 \sum_{j\in J}(\lambda_j/N)^q
     &=N^{-q}\sum_{j\in J}\lambda_j^q,\\
 \frac1N\sum_i(R_i/N)^{q-1}
     &=N^{-q}\sum_i R_i^{q-1}.
\end{align*}
Multiplication by $N^{-q}>0$ and the positive outer power $1/(1-q)$
preserve the inequality, proving the score comparison.
\end{proof}

The analytic step is classical weighted Jensen. Equivalently, restrict
$Z$ to its positive spectral subspace and apply the negative-power
H\"older--McCarthy inequality to the unit vector $Z^{1/2}e_i$
(its squared norm is $Z_{ii}=1$), with exponent $q-1$.
See, for example, the primary treatment in \cite[Theorems~1--2]{Dragomir}.
We do not claim a new general Jensen or operator-power inequality.
The point here is the support-sensitive use of that inequality in the
remaining Vendi--LCR comparison.

\section{The canonical minus-infinity endpoint}
Let $m=\min_{j\in J}\lambda_j>0$ and $r=\min_i R_i>0$.
The canonical endpoint formulas on the same positive support are
\begin{equation}\label{eq:infinity}
 \VS_{-\infty}(Z)=\frac{N}{m},\qquad
 \LCR_{-\infty}(Z)=\frac{N}{r}.
\end{equation}
\begin{theorem}\label{thm:infinity}
For every matrix $Z$ satisfying the stated hypotheses,
$\VS_{-\infty}(Z)\geq\LCR_{-\infty}(Z)$.
\end{theorem}
\begin{proof}
The weights in \eqref{eq:frames} express every $E_i$ as a convex
combination of positive eigenvalues. Hence $m\leq E_i\leq R_i$ for
every $i$, so $0<m\leq r$. Taking reciprocals and multiplying by
$N>0$ proves the claim.
\end{proof}
In particular, neither repeated eigenvalues nor rank one requires a
separate exception. The statement uses the least \emph{positive}
eigenvalue, not the least eigenvalue when zero is present.

\section{Formal coverage and prior-art limits}
The three modules \texttt{VendiJensen}, \texttt{VendiMatrix}, and
\texttt{VendiScores} contain 48 local declarations. A fresh independent
replay compiled every declaration with warnings treated as errors using
Lean~4.34.1, compiler commit
\texttt{5045d0056413266e57c625dcd7c365b10e377c52}, and mathlib revision
\texttt{d13f23b723b8a846827a245b89c10fc7d3f11612} with its pinned
dependencies. Axiom inspection found only \texttt{propext},
\texttt{Classical.choice}, and \texttt{Quot.sound}; there are no proof
holes, custom axioms, or native computation assumptions.

The certificate starts from an actual finite real matrix satisfying the
source hypotheses. It constructs the canonical orthonormal spectral
decomposition, removes only zero-weight terms, proves
\eqref{eq:frames}, derives the energy and row-sum bounds, and proves the
exact normalized comparison. Its endpoints are
\texttt{source\_conjecture\_negative} and
\texttt{source\_conjecture\_negativeInfinity} in namespace
\texttt{VendiScores}. Nonempty dimension and diagonal one prove that the
positive spectrum is nonempty; its positive minimum and the least row
sum are actually selected. No supplied spectral decomposition or
invertibility hypothesis remains. The endpoint \eqref{eq:infinity} is
certified directly; the packet does not include an analytic theorem
about the limit as the order tends to $-\infty$. Known nonnegative-order
results are cited background, not additional formally claimed results.

The source and literature screen was conducted on October~26, 2026.
We read the complete arXiv v4 primary text and its official public
archive, PMC12642762. The later journal PDF required subscription
credentials and returned HTTP~401; we read its public primary abstract
and do not assert that its full body was checked or is byte-identical
to v4. This is a concrete limit on the novelty assessment.

A further prior assertion also deserves acknowledgment. The July~24,
2026 Bioengineering article \cite{Couch}, released in PMC in August,
states the VS--LCR comparison without the rank qualification. Its
supporting reference~50 is the above journal comparison, and its full
public text gives no separate proof of the singular negative-order
case. Thus the general comparison was previously asserted; we claim no
priority for that assertion. The contribution supported here is a
complete proof of the finite support comparison across arbitrary
ranks. Targeted primary-source searches found no earlier demonstrated
resolution of this case or an explicitly stated complete stronger
trace/row-energy bound. This is a bounded prior-art assessment, not an
exhaustive priority certificate. The classical tools, the known cases,
and the equivalent power-sum formulation are not separate results.

\begin{thebibliography}{9}
\bibitem{NguyenPreprint}
P.~Nguyen, J.~Couch, R.~Bansal, A.~Morgan, C.~Tam, M.~Li,
R.~Arnaout and R.~Arnaout,
\emph{Which Similarity-Sensitive Entropy (Sentropy)?},
\href{https://arxiv.org/html/2511.03849v4}{arXiv:2511.03849v4},
March~30, 2026; Theorem~1 and Appendix~C.3.

\bibitem{NguyenJournal}
P.~Nguyen et al.,
\emph{Comparing frameworks for similarity-aware generalized entropy
(sentropy)}, Physical Review E \textbf{113} (2026), 055305,
\href{https://doi.org/10.1103/n2s3-2g85}{doi:10.1103/n2s3-2g85}.
Public abstract inspected; full journal text not accessed.

\bibitem{Dragomir}
S.S.~Dragomir,
\emph{Jensen type weighted inequalities for functions of selfadjoint
and unitary operators}, Italian Journal of Pure and Applied Mathematics
\textbf{32} (2014), 247--264,
\href{https://ijpam.uniud.it/online_issue/201432/23-Dragomir.pdf}{primary PDF}.

\bibitem{Couch}
J.D.~Couch, R.~Arnaout and R.~Arnaout,
\emph{Beyond Size and Class Balance: Alpha as a New Dataset Quality
Metric for Deep Learning}, Bioengineering \textbf{13}(8) (2026), 854,
\href{https://doi.org/10.3390/bioengineering13080854}{doi:10.3390/bioengineering13080854};
public full text \href{https://pmc.ncbi.nlm.nih.gov/articles/PMC13510075/}{PMC13510075}.
\end{thebibliography}
\end{document}
